% from main.tex
%%%%%%%%%%  scan idx 250  |  folio 237  %%%%%%%%%%
\emph{Preuve dans le cas (i).} \(S\) étant localement noethérien, la projection de \(X\) (resp. \(Y\)) sur \(S\) est un morphisme pseudo-cohérent (1.3), et les notions de pseudo-cohérence sur \(X\) (resp. \(Y\)) et relativement à \(S\) coïncident (1.12). D'autre part, comme \(X\) (resp. \(Y\)) est localement noethérien, donc \(\mathcal O_X\) (resp. \(\mathcal O_Y\)) cohérent, un objet \(M\) de \(D(X)\) (resp. \(D(Y)\)) est pseudo-cohérent si et seulement si (I 3.5) les \(H^i(M)\) sont cohérents et \(M\) est localement dans \(D^{-}\). Le théorème résulte donc du suivant :

\textbf{Théorème 2.2.1.} Soit \(f:X\to Y\) un morphisme propre de schémas localement noethériens, et soit \(E\in\Ob D^+(X)\). Si les \(H^i(E)\) sont cohérents, il en est de même des \(R^if_*(E)\). Si en outre \(E\) est localement dans \(D^{-}\), il en est de même de \(\mathbf Rf_*(E)\).

\emph{Preuve.} La première assertion résulte de la suite spectrale
\begin{equation}\tag{$*$}
  R^pf_*(H^q(E)) \Longrightarrow R^*f_*(E)
\end{equation}
et du théorème de finitude (EGA III 3.2.1). Pour la deuxième, on peut supposer \(Y\) affine ; \(X\) est alors quasi-compact et séparé, supposons-le recouvert par \(r\) ouverts affines. On sait (EGA III 1.4.12) que, dans ces conditions, on a \(R^if_*(M)=0\) pour tout \(\mathcal O_X\)-Module quasi-cohérent \(M\) et tout \(i>r\). Si \(E\) est localement dans \(D^{-}\), \(E\) est dans \(D^{-}\) puisque \(X\) est quasi-compact, disons \(H^i(E)=0\) pour \(i>N\). De la suite spectrale (*) on déduit alors que \(R^if_*(E)=0\) pour \(i>N+r\), d'où l'assertion.

\emph{Preuve dans le cas (ii).} La question étant locale sur \(Y\), on peut supposer \(S\) et \(Y\) affines. On a alors un diagramme commutatif à carré cartésien
\[
\begin{tikzcd}[column sep=large,row sep=large]
& & \mathbb P^r_{Y'} \arrow[dr,"p'"] & \\
& \mathbb P^r_Y \arrow[ur,hook,"j"] \arrow[d,"p"'] & & Y' \arrow[d,"g'"]\\
X \arrow[ur,hook,"i"] \arrow[r,"f"] & Y \arrow[urr,hook,"j'"] \arrow[rr,"g"] & & S{,}
\end{tikzcd}
\]
